% Generated draft: requires mathematical and source review.
\documentclass[11pt]{article}
\usepackage[margin=25mm]{geometry}
\usepackage{fontspec}
\setmainfont{Latin Modern Roman}
\usepackage{amsmath,amssymb,mathtools,unicode-math}
\setmathfont{Latin Modern Math}
\usepackage{enumitem,needspace,xurl,hyperref}
\setlength{\parindent}{0pt}
\setlength{\parskip}{5pt}
\setlength{\emergencystretch}{3em}
\newcommand{\N}{\mathbb{N}}
\newcommand{\Z}{\mathbb{Z}}
\newcommand{\Q}{\mathbb{Q}}
\newcommand{\R}{\mathbb{R}}
\newcommand{\C}{\mathbb{C}}
\newcommand{\F}{\mathbb{F}}
\newcommand{\A}{\mathbb{A}}
\newcommand{\PP}{\mathbb{P}}
\newcommand{\E}{\mathbb{E}}
\newcommand{\eps}{\varepsilon}
\newcommand{\dd}{\,\mathrm{d}}
\begin{document}
\section*{UnsolvedMath ID: 700384}
\textbf{Source problem code:} catalogue entry
\subsection{Definitions and mathematical statement}
Work in the usual first-order language of groups, with multiplication, inverse, and identity. A class \(\mathfrak K\) of groups is called \emph{radical} if the following conditions hold:
\begin{enumerate}
\item Whenever \(G\in\mathfrak K\) and \(\varphi:G\twoheadrightarrow H\) is a surjective group homomorphism, then \(H\in\mathfrak K\).
\item Whenever \(G\in\mathfrak K\) and \(N\trianglelefteq G\), then \(N\in\mathfrak K\).
\item For every group \(G\) and every family \((N_i)_{i\in I}\) of normal subgroups of \(G\), if each \(N_i\in\mathfrak K\) and
\[
G=\left\langle\bigcup_{i\in I}N_i\right\rangle,
\]
then \(G\in\mathfrak K\).
\end{enumerate}
A group is \emph{locally finite} if every finitely generated subgroup is finite. A class is locally finite if all its members are locally finite. Write \(\mathrm{Grp}\) for the class of all groups.

The principal question asks whether there is a radical class \(\mathfrak K\ne\mathrm{Grp}\), definable by first-order formulae in the sense intended by the source, which is not locally finite. Thus the required failure of local finiteness is
\[
\exists G\in\mathfrak K\;\exists S\subseteq G\quad
\bigl(S\text{ is finite and }\lvert\langle S\rangle\rvert=\infty\bigr).
\]
The phrase ``definable by formulae'' does not specify whether one sentence or an arbitrary set of sentences is allowed. More explicitly, for a set \(T\) of first-order group sentences, put
\[
\mathfrak K_T=\{G\in\mathrm{Grp}:G\models\sigma\text{ for every }\sigma\in T\}.
\]
Under the theory interpretation, the question asks for such a \(T\) with \(\mathfrak K_T\) radical, proper, and not locally finite. Under the single-sentence interpretation, require \(T=\{\sigma\}\). The supplied record does not resolve which interpretation is intended.

The source also asks the following particular existence question: does there exist a radical class \(\mathfrak K\) such that
\[
\mathbb Z\in\mathfrak K,\qquad
\mathfrak K\ne\mathrm{Grp},\qquad
\prod_{i\in I}G_i\in\mathfrak K
\quad\text{whenever }G_i\in\mathfrak K\text{ for every }i\in I?
\]
Here \(\mathbb Z\) is the infinite cyclic group, and Cartesian products are interpreted as unrestricted direct products indexed by sets \(I\). The particular question does not explicitly impose first-order definability, so that requirement has not been added to it.

As background, the supplied statement asserts that the only non-trivial radical class definable by universal first-order formulae is \(\mathrm{Grp}\), and reports that G. M. Bergman constructed a family of locally finite radical classes definable by first-order formulae, communicated in a letter to the proposer. These are supplied background claims; the record does not define ``non-trivial'', specify the definability convention, or provide proofs or a separate citation for the correspondence.
\subsection{Short English statement}
B. I. Plotkin asks whether a proper radical class of groups can be first-order definable and contain a group that is not locally finite. The meaning of first-order definability requires clarification. A further question asks for a proper radical class containing the infinite cyclic group and closed under Cartesian products; its wording does not explicitly require definability.
\subsection{Sources}
\begin{itemize}
\item Local UnsolvedMath registry, permanent ID \(700384\), source code \texttt{catalogue entry}, title \emph{Kourovka Notebook Problem 6.39}. The supplied registry labels the problem \emph{open}; this label is source data, not an independently verified assessment.
\item \emph{The Kourovka Notebook (Unsolved Problems in Group Theory)}, original problem 6.39, sixth issue (1978), proposed by B. I. Plotkin. The supplied record identifies the twenty-first-issue edition (2026), \texttt{arXiv:1401.0300v48}; its collection description names E. I. Khukhro and V. D. Mazurov as editors.
\item G. M. Bergman's construction is attributed in the supplied problem statement to a letter to the proposer. No separate bibliographic reference, construction, or date for that letter is supplied. The word ``recently'' belongs to the quoted source context and is not treated as a present-day claim.
\item The supplied registry record was created and updated on \texttt{2026-10-16T14:00:00Z}, with upstream revision \texttt{d371ee917ac9647fbb06cd6fc529deec6f3dda82}. These are record timestamps, not literature-verification dates. The cited source and its mathematical background claims have not been independently checked for this draft.
\end{itemize}
\end{document}
